\chapter{Родитель suspension auxiliary copy и Mathai--Quillen достройка}
\label{ch:v10-h15-r2-superconnection-suspension-auxiliary-copy-parent}

Предыдущий аудит выделил superconnection suspension как ближайший источник
второй копии $\Sigma$. Проверим различие между градуированным удвоением и
независимой бозонной координатой, а затем минимально достроим suspension до
Mathai--Quillen Thom-пары.

На удвоенном носителе $\Sigma\oplus\Sigma$ положим
\begin{equation}
 \Gamma_{\rm sus}=\begin{pmatrix}I_8&0\\0&-I_8\end{pmatrix},
 \qquad
 S_{\rm sus}=\begin{pmatrix}0&I_8\\I_8&0\end{pmatrix}.
\end{equation}
Тогда $S_{\rm sus}^2=I_{16}$, $\{\Gamma_{\rm sus},S_{\rm sus}\}=0$, а
$S_{\rm sus}$ коммутирует с удвоенными gauge- и Real-действиями. Тем самым
suspension точно создаёт копию, grading и odd edge. Однако отображение
независимых динамических полей текущего parent во второе слагаемое имеет
rank ноль. Само удвоение является стабилизацией, а не происхождением
$A_\Sigma$.

Минимальная Thom-достройка вводит нечётную координату $\chi_\Sigma$ и
бозонную координату $H_\Sigma$ того же восьмимерного bundle-типа. Их
линейный contractible differential имеет вид
\begin{equation}
 d_{\rm MQ}=\begin{pmatrix}0&0\\I_8&0\end{pmatrix},
 \qquad d_{\rm MQ}^2=0,
 \qquad \rank d_{\rm MQ}=8.
\end{equation}
Поэтому kernel и image имеют dimension восемь, а cohomology пары равна
нулю: $H_\Sigma$ не является новой on-shell частицей. Его injection имеет
rank восемь и после действия $P_{\rm rel}$ возбуждает все восемь relative
мод.

Уже выведенная mixed curvature задаёт Thom-section
\begin{equation}
 s_Q(\Sigma)=Q\Sigma,
 \qquad
 Q=\operatorname{diag}(3,-3,7,1,-1,-7,3,-3).
\end{equation}
Для unit bundle metric совместный квадратичный блок $(H_\Sigma,\Sigma)$
равен
\begin{equation}
 \mathcal H_{\rm MQ}=
 \begin{pmatrix}I_8&Q\\Q&49I_8\end{pmatrix}.
\end{equation}
Точное Gaussian elimination даёт Schur complement
\begin{equation}
 49I_8-Q^2
 =\operatorname{diag}(40,40,0,48,48,0,40,40),
\end{equation}
имеющий inertia $(0_-,2_0,6_+)$. Линейный fermionic Jacobian секции
постоянен и равен $\det Q=3969$. При фиксированном $Q$ он может войти в
нормировку меры, но при динамическом $\Delta$-фоне требуется полный
curvature/Jacobian audit.

\begin{theorem}[Suspension не равна auxiliary origin]
Чистая superconnection suspension закрывает только $3/7$ критериев: она
даёт типизированное положительное градуированное удвоение без новой
on-shell моды, но не даёт независимой переменной, nilpotent pair или
section $Q\Sigma$. Mathai--Quillen-достройка условно закрывает $6/7$:
остаётся вывести её cohomological differential и полную меру из текущего
parent.
\end{theorem}

\begin{nogocriterion}[Стабилизацию нельзя считать новым полем]
Наличие второго изоморфного слагаемого и odd swap-оператора не доказывает
существование независимо интегрируемого $A_\Sigma$. Без contractible пары,
section и меры такое удвоение является только алгебраической
стабилизацией.
\end{nogocriterion}

\begin{protocol}\begin{enumerate}
\item suspension carrier dimension: \textbf{$16$};
\item suspension edge rank: \textbf{$16$};
\item bare dynamical injection rank: \textbf{$0$};
\item Thom differential rank/cohomology: \textbf{$8/0$};
\item independent Thom boson rank: \textbf{$8$};
\item relative image rank: \textbf{$8$};
\item section rank and Jacobian: \textbf{$8$, $3969$};
\item Schur inertia: \textbf{$(0_-,2_0,6_+)$};
\item bare/conditional status: \textbf{$3/7$ и $6/7$};
\item следующий гейт: \textbf{общий parent Mathai--Quillen Thom-мультиплета}.
\end{enumerate}\end{protocol}

Литературные основания: Quillen \cite{v10-sus-quillen},
Mathai--Quillen \cite{v10-sus-mathai-quillen} и полевой разбор Blau
\cite{v10-sus-blau}.

Машинный сертификат сохранён в
\path{s2t_v10_particle_wrinkle_dislocation_callias_equal_charge_twist_m4_h15_r2_superconnection_suspension_auxiliary_copy_parent_origin_gate_results.json}.

\begin{thebibliography}{9}
\bibitem{v10-sus-quillen}
D.~Quillen, \emph{Superconnections and the Chern Character}, Topology 24
(1985), 89--95.
\bibitem{v10-sus-mathai-quillen}
V.~Mathai, D.~Quillen, \emph{Superconnections, Thom Classes, and
Equivariant Differential Forms}, Topology 25 (1986), 85--110.
\bibitem{v10-sus-blau}
M.~Blau, \emph{The Mathai--Quillen Formalism and Topological Field Theory},
J. Geom. Phys. 11 (1993), 95--127, arXiv:hep-th/9203026.
\end{thebibliography}